\documentclass[11pt]{article}
\usepackage[margin=1in]{geometry}
\usepackage{amsmath,amssymb,booktabs,array}
\usepackage[hidelinks]{hyperref}

\title{Bilal Search Slides and the Low-Rank Modeling Section}
\author{}
\date{\today}

\newcommand{\E}{\mathbb{E}}
\newcommand{\1}{\mathbf{1}}

\begin{document}
\maketitle

\section*{Source Convention}

Page citations to Bilal refer to the PDF pages in
\texttt{bilal\_2024\_slides\_search-and-matching\_Lecture2.pdf}. Page citations
to Bonhomme--Lamadon--Manresa refer to the printed Econometrica pages in BLM
(2019). Page citations to the project note refer to PDF pages in
\texttt{loo\_full\_note\_v1.pdf}. When I reuse results from the prior project
notes, I cite the prior result PDF \texttt{7\_15\_Meeting (1).pdf} and the
source \TeX{} file \texttt{shimer\_smith\_writeup.tex}. Displays or claims not
directly sourced are labeled as derived synthesis, derived ANOVA steps, or
derived nesting interpretations.

\section{Notation}

\begin{center}
\begin{tabular}{>{\raggedright\arraybackslash}p{0.24\linewidth}
                >{\raggedright\arraybackslash}p{0.68\linewidth}}
\toprule
Symbol & Meaning \\
\midrule
$u_t,v_t$ & Bilal unemployment and vacancy stocks at time $t$. Source: Bilal, pp. 9--10. \\
$M_t,M(u,v)$ & Bilal aggregate matching function and number of matches. Source: Bilal, p. 9. \\
$m,\alpha$ & Bilal Cobb--Douglas matching-function scale and unemployment elasticity in $M(u,v)=mu^\alpha v^{1-\alpha}$. Source: Bilal, p. 9. \\
$\theta_t$ & Bilal labor-market tightness, $v_t/u_t$. Source: Bilal, p. 10. \\
$f_t,q_t$ & Bilal job-finding and vacancy-filling probabilities. Source: Bilal, p. 10. \\
$s$ & Bilal exogenous job-termination probability. Source: Bilal, p. 10. \\
$U_t,E_t,V_t,J_t$ & Bilal values of unemployment, employment, a vacant job/firm, and a filled job/firm. Source: Bilal, p. 14. \\
$b,c,z_t,w_t,\beta$ & Bilal unemployment flow payoff, vacancy cost, productivity, wage, and discount factor. Source: Bilal, pp. 9, 14. \\
$S_t,\gamma,r$ & Bilal joint surplus, worker bargaining power, and steady-state interest-rate notation $r=\beta^{-1}-1$. Source: Bilal, pp. 15, 19, 23. \\
$Y_{ij},X_{ij}'\delta,\varepsilon_{ij}$ & Project wage/log-wage observation, observed covariate component, and residual in the low-rank interface. Source: project note, p. 8. \\
$i,j,I,J$ & Worker and firm indices and counts in the project interface. Source: project note, p. 8. \\
$\alpha_i,\psi_j,u_i,v_j,\Lambda,\operatorname{rank}(\Lambda)=r$ & Project additive worker and firm components, worker and firm interaction factors, interaction matrix, and interaction rank. Source: project note, p. 8. The symbol $r$ is therefore context-dependent: an interest rate in Bilal's steady-state slide and an interaction rank in the project interface. \\
$\delta,\varepsilon_{ij}$ & Project covariate coefficient vector in $X_{ij}'\delta$ and residual in $(\star)$. Source: project note, p. 8. \\
$k,\ell,K,L$ & Worker type and firm class indices and counts in the saturated type-cell construction. Source: project note, p. 8. \\
$\mu_{k\ell},\mu_{\cdot\cdot},\mu_{k\cdot},\mu_{\cdot\ell}$ & Saturated type-cell mean and its grand, row, and column means. Derived from the ANOVA decomposition in project note Proposition 1, p. 8. \\
$\bar\alpha_k,\bar\psi_\ell,M_{k\ell}$ & Additive type components and double-centered interaction matrix in the ANOVA decomposition. Derived from project note Proposition 1, p. 8. This $M$ is unrelated to Bilal's matching function $M(u,v)$. \\
$e_k,e_\ell,\1_K,\1_L$ & Standard basis vectors and all-ones vectors used in the type-indicator nesting and rank proof. Derived notation for project note Proposition 1, p. 8. \\
$k_{it},\ell(j),\alpha_i,a_t(k),b_t(k),X_{it}'c_t,\varepsilon_{it}$ & BLM firm class, project firm-class map, worker type, class intercept, class slope, covariate component, and residual in the simple static regression example. Source for BLM objects: BLM, pp. 702--704; project notation $\ell(j)$ is derived for the type-cell embedding. \\
\bottomrule
\end{tabular}
\end{center}

\section{Bilal Lecture 2: What Matters for This Project}

Bilal's Lecture 2 is a compact DMP search-and-matching benchmark. It is not a
two-sided worker-type/firm-type model like Shimer--Smith or
Borovi\v{c}kov\'a--Shimer, but it gives the reusable search logic: matching
technology, value functions, joint surplus, Nash bargaining, free entry, and
the unemployment stock-flow law.

\subsection{Matching Technology and Tightness}

The lecture defines vacancies $v_t$ and unemployment $u_t$, with matches formed
by a matching function
\[
M_t=M(u_t,v_t).
\]
The slide specializes to Cobb--Douglas,
\[
M(u,v)=m u^\alpha v^{1-\alpha},
\]
while noting that the qualitative results extend to any constant-returns
matching function. Source: Bilal, p. 9.

Labor-market tightness is
\[
\theta_t=\frac{v_t}{u_t}.
\]
The job-finding probability for an unemployed worker is
\[
f_t=\frac{M_t}{u_t}=f(\theta_t)=m\theta_t^{1-\alpha},
\]
and the vacancy-filling probability is
\[
q_t=\frac{M_t}{v_t}=q(\theta_t)=m\theta_t^{-\alpha}
=\frac{f(\theta_t)}{\theta_t}.
\]
Source: Bilal, p. 10. The key conceptual point for our work is that random
matching can be summarized either at the individual meeting-rate level, as in
SS and BS, or at the aggregate matching-function/tightness level, as in DMP.

\subsection{Value Functions}

The lecture's four Bellman equations are
\[
U_t=b+\beta \E_t\left[f_t E_{t+1}+(1-f_t)U_{t+1}\right],
\]
\[
E_t=w_t+\beta \E_t\left[sU_{t+1}+(1-s)E_{t+1}\right],
\]
\[
V_t=-c+\beta \E_t\left[q_tJ_{t+1}+(1-q_t)V_{t+1}\right],
\]
\[
J_t=z_t-w_t+\beta \E_t\left[sV_{t+1}+(1-s)J_{t+1}\right].
\]
Source: Bilal, p. 14. This is the slide-deck analogue of the Bellman logic used
in the SS and BS memo: current payoff plus expected continuation values, with
transitions generated by search and separation.

\subsection{Surplus and the Wage}

The worker surplus is $E_t-U_t$, the firm surplus is $J_t-V_t$, and free entry
sets $V_t=0$. The joint match surplus is
\[
S_t=E_t+J_t-U_t,
\]
and the lecture derives
\[
S_t
=
z_t-b-\beta f_t\E_t[E_{t+1}-U_{t+1}]
+\beta(1-s)\E_t[S_{t+1}].
\]
Source: Bilal, p. 15; the negative sign on the foregone-search term is also
confirmed by the later joint-surplus equation on Bilal, p. 21,
$S_t=z_t-b+\beta(1-s-\gamma f_t)\E_t[S_{t+1}]$. The useful point is the same
one that makes SS and BS clean: with transferable utility and linear payoffs,
joint surplus is independent of the wage. The wage only splits the surplus.

Bilal then imposes generalized Nash bargaining:
\[
w_t=\arg\max_w (E_t(w)-U_t)^\gamma J_t(w)^{1-\gamma},
\]
so
\[
E_t(w_t)-U_t=\gamma S_t,
\qquad
J_t(w_t)=(1-\gamma)S_t.
\]
Source: Bilal, p. 19. The equilibrium wage can be written as
\[
w_t=(1-\gamma)b+\gamma(z_t+\theta_t c),
\]
after using free entry. Source: Bilal, p. 20. For our project, this is a
warning and a guide: search frictions determine outside options and surplus,
but the wage equation depends on the bargaining/transfer rule. That is why the
SS-to-BS mapping had to track outside values and surplus shares carefully.

\subsection{Equilibrium Conditions}

Given surplus sharing, the joint-surplus Bellman becomes
\[
S_t=z_t-b+\beta(1-s-\gamma f_t)\E_t[S_{t+1}],
\]
free entry determines tightness,
\[
c=\beta q_t(1-\gamma)\E_t[S_{t+1}],
\]
and unemployment follows
\[
u_{t+1}-u_t=s(1-u_t)-f_tu_t.
\]
Source: Bilal, p. 21. In steady state, these reduce to the job-creation condition
\[
\frac{r+s+\gamma f(\theta)}{(1-\gamma)q(\theta)}
=
\frac{z-b}{c},
\qquad r=\beta^{-1}-1,
\]
and the Beveridge-curve stock-flow identity
\[
u=\frac{s}{s+f(\theta)}.
\]
Source: Bilal, p. 23.

\subsection{Connection to SS and BS}

Derived synthesis: Bilal's slides should be used as the generic search-model scaffold, not as a
direct nesting target. SS and BS add type-pair heterogeneity to this search
logic. SS makes match output deterministic conditional on types, $f(x,y)$. BS
adds a match-specific draw $z$ and therefore a type-pair selection threshold.
The shared structure is:
\[
\text{Bellman equations}
\quad\Rightarrow\quad
\text{joint surplus independent of wage}
\quad\Rightarrow\quad
\text{Nash split}
\quad\Rightarrow\quad
\text{acceptance/job-creation condition}.
\]

\section{Modeling Section: The Low-Rank Interface}

The project note's modeling object is
\[
Y_{ij}
=
X_{ij}'\delta+\alpha_i+\psi_j+u_i^\top\Lambda v_j+\varepsilon_{ij},
\qquad
\operatorname{rank}(\Lambda)=r\ll \min\{I,J\}.
\tag{$\star$}
\]
Source: project note, p. 8; prior result PDF, pp. 18--21;
\texttt{shimer\_smith\_writeup.tex}, Sec. 6. The point of the modeling section is
not to claim that every structural model is literally inside $(\star)$. The point
is to identify which static wage-equation cores land in this interface after
normalization.

\subsection{The Double-Demeaned $K\times L$ Construction}

Let worker types be $k\in\{1,\ldots,K\}$ and firm classes be
$\ell\in\{1,\ldots,L\}$. Let
\[
\mu=(\mu_{k\ell})\in\mathbb R^{K\times L}
\]
be the saturated static type-pair mean wage surface:
\[
\mu_{k\ell}
=
\E[Y\mid \text{worker type }k,\text{ firm class }\ell].
\]
This is the object shown in the Saturday screenshot
\texttt{C:\textbackslash Users\textbackslash A8327\textbackslash Downloads\textbackslash image (3).png}
and in the project note's Proposition 1. Source: project note, p. 8.

Derived ANOVA step: define the grand mean, row means, and column means by
\[
\mu_{\cdot\cdot}
=
\frac{1}{KL}\sum_{k=1}^K\sum_{\ell=1}^L\mu_{k\ell},
\qquad
\mu_{k\cdot}
=
\frac{1}{L}\sum_{\ell=1}^L\mu_{k\ell},
\qquad
\mu_{\cdot\ell}
=
\frac{1}{K}\sum_{k=1}^K\mu_{k\ell}.
\]
Set
\[
\bar\alpha_k=\mu_{k\cdot}-\mu_{\cdot\cdot},
\qquad
\bar\psi_\ell=\mu_{\cdot\ell},
\]
and define the double-centered interaction matrix
\[
M_{k\ell}
=
\mu_{k\ell}-\mu_{k\cdot}-\mu_{\cdot\ell}+\mu_{\cdot\cdot}.
\]
Then
\[
\mu_{k\ell}
=
\bar\alpha_k+\bar\psi_\ell+M_{k\ell}.
\]
This is a derived ANOVA decomposition. The interaction matrix is double-centered:
\[
\sum_{\ell=1}^L M_{k\ell}=0\quad\text{for every }k,
\qquad
\sum_{k=1}^K M_{k\ell}=0\quad\text{for every }\ell.
\]
It is unique once the mean convention above is fixed.

Derived embedding in $(\star)$: assign worker $i$ to type $k(i)$ and firm $j$ to class $\ell(j)$. Let
$e_k\in\mathbb R^K$ and $e_\ell\in\mathbb R^L$ be standard basis vectors. Set
\[
\alpha_i=\bar\alpha_{k(i)},
\qquad
\psi_j=\bar\psi_{\ell(j)},
\qquad
u_i=e_{k(i)},
\qquad
v_j=e_{\ell(j)},
\qquad
\Lambda=M.
\]
Then
\[
u_i^\top\Lambda v_j
=
e_{k(i)}^\top M e_{\ell(j)}
=
M_{k(i)\ell(j)},
\]
and therefore
\[
\alpha_i+\psi_j+u_i^\top\Lambda v_j
=
\mu_{k(i)\ell(j)}.
\]
So the saturated $K\times L$ type-pair mean surface is exactly a special case of
$(\star)$.

\subsection{Rank Bound}

Derived rank bound: because $M$ is double-centered, its columns lie in the $K-1$ dimensional
subspace orthogonal to $\1_K$, and its rows lie in the $L-1$ dimensional
subspace orthogonal to $\1_L$. Hence
\[
\operatorname{rank}(M)
\leq
\min\{K-1,L-1\}
=
\min\{K,L\}-1.
\]
This is the rank bound in the screenshot and in the project note. Source:
project note, p. 8. The important interpretation is that the additive row and
column components remove one dimension from each side. The interaction matrix
does not live in all of $\mathbb R^{K\times L}$; it lives between centered type
spaces.

\section{In What Sense We Nest BLM}

BLM's static model uses firm classes and worker heterogeneity. Firms are grouped
into classes $k(j)$, and workers have latent types $\alpha_i$, which may be
discrete or continuous. BLM are interested in earnings distributions for workers
of a given type in a given firm class, as well as type proportions and sorting.
Source: BLM, pp. 702--704. A simple static BLM regression example is
\[
Y_{it}=a_t(k_{it})+b_t(k_{it})\alpha_i+X_{it}'c_t+\varepsilon_{it}.
\]
Source: BLM, p. 704.

Derived nesting interpretation: the honest nesting claim has three layers.

\paragraph{1. Exact for static type-cell mean surfaces.}
If BLM worker heterogeneity is discretized into $K$ worker types and firm
heterogeneity into $L$ firm classes, then the static conditional mean surface
\[
\mu_{k\ell}
=
\E[Y\mid k(i)=k,\ell(j)=\ell]
\]
is exactly nested by $(\star)$ through the double-demeaned construction above.
This is the clean reading of the screenshot's BLM row.

\paragraph{2. Exact for BLM's simple interactive mean example.}
The displayed BLM regression
\[
a_t(k)+b_t(k)\alpha_i
\]
is a firm-class intercept plus a firm-class slope on worker type. For a fixed
date $t$, it is nested by
\[
\psi_j=a_t(\ell(j)),
\qquad
u_i=\alpha_i,
\qquad
v_j=b_t(\ell(j)),
\qquad
\Lambda=1,
\]
with observed covariates absorbed into $X_{ij}'\delta$. This is rank one if
$\alpha_i$ is scalar. If the empirical object is instead a fully saturated
discrete type-class mean matrix, the exact representation is the $M$ construction
above, with rank at most $\min\{K,L\}-1$.

\paragraph{3. Not exact for the full BLM framework.}
The full BLM contribution is larger than a static mean equation. It includes
latent firm-class estimation, distributional worker heterogeneity, mixture
identification, mobility assumptions, and a dynamic model where earnings can
affect mobility and post-move earnings can depend on previous earnings and
previous firm class. Source: BLM, pp. 701--706. Those layers are not themselves
nested by $(\star)$, because $(\star)$ is a static wage-interface equation, not a
full distributional panel estimator or dynamic mobility model.

So the phrase ``we nest BLM'' is correct only if it is expanded as:
\[
\text{we nest BLM's static wage-equation/type-cell mean layer.}
\]
It should not be written as:
\[
\text{we nest the full BLM empirical framework.}
\]
This matches the prior result PDF's warning that the factor interface nests
static wage-equation cores, while BLM's grouping, mixture estimation, and dynamic
extensions remain separate implementation and interpretation layers. Source:
prior result PDF, pp. 18--21; \texttt{shimer\_smith\_writeup.tex}, Sec. 6.5.

\section{Clean Wording for the Modeling Section}

The safest paragraph is:

\begin{quote}
The low-rank interface nests the static wage-equation cores of several models.
For BLM, the exact statement is about the static type-cell mean surface. If
worker types are $k=1,\ldots,K$ and firm classes are $\ell=1,\ldots,L$, any
saturated mean matrix $\mu_{k\ell}$ decomposes uniquely into additive row and
column terms plus a double-centered interaction matrix $M$. Taking worker and
firm factors to be type indicators and $\Lambda=M$ reproduces the mean surface
exactly, with $\operatorname{rank}(M)\leq\min\{K,L\}-1$. This nests BLM's static
mean layer, but not BLM's full distributional estimator, class-estimation step,
or dynamic panel model.
\end{quote}

\end{document}
